\documentclass[sigconf,nonacm]{acmart}
\settopmatter{printacmref=false}
\setcopyright{none}
\renewcommand\footnotetextcopyrightpermission[1]{}
\pagestyle{plain}
\raggedbottom

\acmConference[Manuscript Draft]{Count-Layer Audit Study}{2026}{}
\acmYear{2026}
\acmISBN{}
\acmDOI{}

\title{Count-Layer Drift in Nwagu Aneke: Separating Source Evidence, Derived Orthographic Expansion, and Public Claims}
\author{Anonymous Author(s)}
\affiliation{\institution{Independent research manuscript}\country{Nigeria}}
\email{anonymous@example.invalid}

\begin{abstract}
Underdocumented scripts are often represented through tables, teaching charts,
transcriptions, source images, and later analytical normalizations. These
objects can be useful at the same time that they support different kinds of
claims. This paper examines the problem through Nwagu Aneke, an Igbo writing
system whose current working evidence base contains a 26 by 8 source-observed
table, yielding 208 records, and a 27/216 representation produced only when a
combined f/v row is split for analysis. The paper defines count-layer drift as
the error in which a number from one evidentiary layer is repeated as though it
belonged to another layer. A six-field audit model records quantity, layer,
source basis, transformation rule, permitted wording, and revision condition.
Applied to the Nwagu Aneke case, the model preserves 26 by 8 as the present
source-observed count and treats 27/216 as a derived orthographic expansion.
The result is a reproducible source-critical boundary for digital-humanities,
standards-readiness, and computational representations of the script: later
systems may use a derived f/v view, but they must not report it as the observed
source inventory.
\end{abstract}

\keywords{Nwagu Aneke, Igbo writing systems, source criticism, digital humanities, evidence layers, script documentation}

\begin{document}
\maketitle

\section{Introduction}
Digital scholarship on underdocumented scripts depends on acts of
representation. A manuscript image becomes a crop, a chart becomes a table, a
row label becomes metadata, a classroom form becomes a teaching interface, and
an analytical normalization becomes a computational inventory. Each act can be
useful, but each act can also change the evidentiary status of the object being
represented. When the change is not recorded, a count can drift. A number that
belongs to a source table can be repeated as a character repertoire; a
normalization introduced for analysis can be cited as a historical observation;
an unresolved chart can become a software inventory.

This paper studies that failure mode through Nwagu Aneke, an Igbo writing
system discussed in scholarship and public script-status resources
\cite{azuonye1992,omniglotNwagu,scriptSourceNwagu,sirisNwaguProposal}. The
current working source layer used here is a 26 by 8 table of row and
vowel/modifier combinations, for 208 source-observed records. A second number,
27/216, appears only when a combined f/v row is split as a derived analytical
operation. Both representations can be useful. They cannot, however, support the
same claim. The 26 by 8 table describes the present source-observed structure
used by this study. The 27/216 representation describes a transformation applied
to that structure. If the transformed representation is later described as the
source-observed inventory, a derived model has been promoted into a source
claim.

The research question is: can a small count-layer audit prevent such promotion
while still allowing derived representations to be used responsibly? The
hypothesis is that count claims become more reliable when every quantity is
stored together with its layer, source basis, transformation rule, permitted
wording, and revision condition. Under that model, the source-observed
26 by 8 layer and the derived 27/216 f/v expansion can coexist without being
collapsed into a single inventory claim.

The paper's central result is a source-critical decision rule. In the current
evidence base, 26 by 8 equals 208 is the source-observed layer. The 27/216
representation is permitted only as a derived f/v expansion. A manuscript,
dataset, teaching tool, tokenizer, standards note, or visualization that uses
27/216 must label it as derived. A public statement that treats 27/216 as the
observed source inventory fails the audit. This result does not settle the full
history of Nwagu Aneke, produce a public glyph corpus, or establish a Unicode
repertoire. It supplies a narrower but important control: which count may be
used for which kind of claim under the evidence now available.

The contribution matters beyond this single case because count-layer drift is a
general risk in cultural-heritage computing. Digital-humanities standards can
represent source witnesses, annotations, entities, derivations, and research
objects \cite{teiP5Glyphs,iiifPresentation3,w3cAnnotationModel,w3cAnnotationVocab,w3cProvO,rocrate12,cidocCrm,datacite45}.
FAIR and CARE principles clarify that preservation, reuse, and authority are
not the same problem \cite{wilkinson2016fair,carroll2020care}. Unicode guidance
distinguishes script interest from the evidentiary requirements for encoding
\cite{unicodeAfricanScripts2023,unicode17CoreSpec,unicodeProposals,unicodeRoadmaps,unicodePipeline2026,seiProposalTips}.
Tokenization and claim-verification research show that representational units
and fluent summaries can affect downstream conclusions
\cite{sennrich2016bpe,kudo2018sentencepiece,bostrom2020bytepair,rust2021tokenization,wadden2020scifact,gao2023alce}.
What these literatures do not automatically supply is the local source-critical
judgment that tells a project whether a particular count is observed, derived,
speculative, blocked, or revised.

\section{Related Work}
The source anchor for this study is Azuonye's account of the Nwagu Aneke script,
together with public script-status resources that identify the script and its
relationship to broader discussions of African writing systems
\cite{azuonye1992,omniglotNwagu,scriptSourceNwagu,sirisNwaguProposal}. These
sources are sufficient to motivate a source-critical and digital-humanities
study. They are not, by themselves, sufficient to settle a public character
repertoire, a representative glyph set, or an authority decision. That
distinction is precisely why a count-layer audit is needed.

Digital-humanities infrastructure provides mature tools for describing cultural
materials while preserving the fact that some representations are provisional.
TEI can record
characters, glyphs, uncertain readings, and editorial practice
\cite{teiP5Glyphs}. IIIF can identify canvases and regions in image-based
collections \cite{iiifPresentation3}. The Web Annotation Data Model and
Vocabulary can attach assertions, comments, and classifications to source
targets \cite{w3cAnnotationModel,w3cAnnotationVocab}. PROV-O can represent
entities, activities, agents, and derivations \cite{w3cProvO}. RO-Crate and
DataCite support research-object packaging and citation
\cite{rocrate12,datacite45}. CIDOC CRM provides a cultural-heritage ontology for
objects, events, actors, and interpretations \cite{cidocCrm}. These standards
are important because they can carry the result of a count-layer audit. They do
not perform the audit automatically.

Responsible data governance is also relevant. FAIR principles emphasize
findability, accessibility, interoperability, and reuse
\cite{wilkinson2016fair}. CARE principles emphasize collective benefit,
authority to control, responsibility, and ethics \cite{carroll2020care}. For an
Igbo writing-system artifact, those two framings must be read together. A table
may be technically reusable while still requiring source review, rights review,
and community authority review before images, examples, or representative forms
are circulated. The count-layer audit therefore treats rights and authority as
part of the evidentiary state, not as a decorative limitation at the end of the
paper.

Standards-facing work adds a second layer of caution. Unicode documentation and
Script Encoding Working Group guidance distinguish scripts, characters, glyphs,
properties, examples, names, fonts, proposal documents, and roadmaps
\cite{unicodeAfricanScripts2023,unicode17CoreSpec,unicodeProposals,unicodeRoadmaps,unicodePipeline2026,seiProposalTips}.
An underdocumented script can be important, named, and worth preserving while
still lacking the evidence required for a formal encoding proposal. The
Nwagu Aneke case illustrates why. A 26 by 8 source table can support a
readiness audit and source-critical discussion, but it cannot be treated as a
completed proposal dossier. A derived f/v expansion can support modeling, but it
cannot stand in for observed source evidence.

Computational representation research clarifies the downstream stakes. Subword
tokenization work shows that unit choices affect language models and later
systems \cite{sennrich2016bpe,kudo2018sentencepiece,bostrom2020bytepair,rust2021tokenization}.
Scientific claim-verification and citation-grounded generation research show
that text can be fluent while misrepresenting what evidence supports
\cite{wadden2020scifact,gao2023alce}. The same risk appears in script
documentation. A table that has been normalized for analysis can be rendered
cleanly, cited smoothly, and used in software while no longer preserving the
source layer from which it came. Count-layer auditing is one way to make that
representational shift visible.

\section{Materials and Evidence}
The material examined in this paper is a text-based source and analysis layer
for Nwagu Aneke. The source-observed structure is treated as 26 rows by
8 vowel/modifier columns, yielding 208 records. The f/v expansion is a derived
operation in which a combined f/v row is separated into two analytical rows,
producing 27 rows and 216 cells. The difference between those numbers is not a
minor notation issue. It marks the difference between a source-observed table
and a transformation applied to that table.

The paper uses three kinds of evidence. The first is source evidence: the local
record of the 26 by 8 table and the scholarly or public sources that identify
Nwagu Aneke. The second is transformation evidence: the operation that produces
the 27/216 representation by splitting f/v. The third is claim-boundary
evidence: standards, provenance, annotation, and responsible-data literature
that explains why a representation must preserve its evidence layer. The paper
does not depend on public reproduction of manuscript images or glyph crops. It
does not require an implementation artifact such as a font or keyboard. Those
would be separate evidentiary objects.

Table~\ref{tab:count_layers} summarizes the count layers examined here. The
table is deliberately simple because the main issue is not arithmetic. The main
issue is whether the arithmetic is attached to the correct kind of claim.

\begin{table}[t]
\caption{Count layers used in the audit}
\label{tab:count_layers}
\begin{tabular}{p{0.28\columnwidth}p{0.22\columnwidth}p{0.40\columnwidth}}
\toprule
Layer & Quantity & Permitted interpretation \\
\midrule
Source-observed table & 26 by 8 = 208 & Current table of row and vowel/modifier records \\
Derived f/v expansion & 27 by 8 = 216 & Analytical view produced by splitting a combined f/v row \\
Standards repertoire & unresolved & Requires separate character, glyph, usage, and authority evidence \\
Public data release & unresolved & Requires rights and source-access decisions \\
\bottomrule
\end{tabular}
\end{table}

Two boundaries follow from the materials. First, a derived count can be useful
without being source-observed. The 27/216 representation may help compare
orthographic patterns, build teaching views, test layer-preserving software, or
construct a normalized analytical table. Second, usefulness is not the same as
evidentiary status. A source-critical paper must not cite the derived count as
the observed inventory simply because it is tidier or more convenient for a
downstream task.

\section{Method}
The method is a six-field count-layer audit. Each count claim is represented by
the following fields: quantity, evidence layer, source basis, transformation
rule, permitted wording, and revision condition. A count is acceptable only when
all six fields can be stated. If a count has a quantity but no layer, it is
ambiguous. If it has a layer but no source basis, it is unsupported. If it has a
transformation rule but the wording hides that transformation, it risks drift.
If it has no revision condition, later evidence cannot correct it cleanly.

The evidence layers used here are source-observed, derived, standards-facing,
speculative, and blocked. Source-observed counts are those supported by the
present table or source record under discussion. Derived counts are
transformations made from source-observed material. Standards-facing counts are
candidate units for character encoding or interchange and require a different
evidentiary package. Speculative counts may motivate future work but cannot be
reported as results. Blocked counts are those for which the relevant evidence
cannot yet be inspected, released, or authorized.

The audit proceeds in four passes. The first pass identifies every count in the
manuscript or dataset. In this case the relevant counts are 26, 8, 208, 27, and
216. The second pass assigns each count to a layer. The third pass writes
permitted and forbidden wording. For example, ``the current source-observed
table has 26 rows by 8 columns'' is permitted under the present evidence base;
``the script has 216 source-observed signs'' is not. The fourth pass states the
revision condition that would change the result. A reviewed source record could
alter the source table. A separate standards analysis could propose different
characters. A rights decision could allow public examples. Each change would
affect a specific layer rather than silently rewriting all layers at once.

The method includes negative controls. The primary negative control is any
sentence that describes 27/216 as source-observed. A second negative control is
any table caption, abstract sentence, or conclusion sentence that reports only
the derived number without naming its transformation. A third negative control
is any standards-facing statement that treats the table as a completed Unicode
repertoire. A fourth negative control is any public-data statement that implies
source images or glyph crops may be released without rights and authority
review. These controls are intentionally linguistic because drift often appears
first in prose before it appears in data structures.

The method can be represented as a rule:

\begin{quote}
A count may be reused only if its evidence layer and transformation rule remain
attached to the statement in which the count appears.
\end{quote}

The rule is stricter than ordinary citation checking. A sentence can cite a real
source and still promote the wrong layer. It is also stricter than ordinary
format checking. A formatted table can be beautiful and still wrong if it
collapses source and derived counts. The audit therefore asks whether the
claim, not just the citation or format, has preserved the evidence layer.

\section{Results}
The audit produces three results. First, the present source-observed count is
26 by 8, yielding 208 records. This statement is a source-layer statement. It is
the appropriate basis for source-critical discussion under the current evidence
base.

Second, the 27/216 representation is derived. It is produced when the combined
f/v row is split into separate analytical rows. The f/v split may be useful for
orthographic modeling, teaching, search, comparison, or later computational
experiments, but it does not change the source-observed layer. The derived
representation must therefore be written as derived wherever it appears.

Third, the audit blocks three stronger moves. It blocks a source claim that
reports 27/216 as observed. It blocks a standards claim that treats either count
as a completed character repertoire. It blocks a public-data claim that implies
source images, manuscript examples, representative glyphs, or community
authority have been cleared for release. These blocks are not failures of the
audit. They are the audit's main value: they prevent a useful representation
from becoming a false foundation.

Table~\ref{tab:permitted_wording} gives examples of the wording distinction.
The left column contains formulations that preserve the layer. The right column
contains formulations that fail the audit.

\begin{table*}[t]
\caption{Permitted and failed wording under the count-layer audit}
\label{tab:permitted_wording}
\begin{tabular}{p{0.46\textwidth}p{0.46\textwidth}}
\toprule
Permitted wording & Failed wording \\
\midrule
The current source-observed table has 26 rows by 8 vowel/modifier columns. &
Nwagu Aneke has a 216-sign source inventory. \\
The 27/216 representation is produced by a derived f/v split. &
The source table contains 27 rows and 216 observed cells. \\
The f/v split may be used as a labeled analytical view. &
The f/v split proves a completed repertoire. \\
The evidence supports a source-critical count audit. &
The evidence supplies a Unicode proposal or public glyph corpus. \\
\bottomrule
\end{tabular}
\end{table*}

The result is narrow but consequential. It changes what later work may safely
say. A digital edition can show the 26 by 8 table as the present source layer
and offer a derived f/v view as a separate mode. A teaching interface can use
the split only if learners can still see that it is derived. A tokenizer
experiment can use the derived view only if its evaluation reports the
transformation. A readiness matrix can cite the 26 by 8 table as evidence for
discussion while reserving repertoire, glyph, name, and implementation claims
for later work.

\section{The f/v Hinge as a Case Study}
The f/v hinge is the most important derived operation in the present case
because it produces the cleanest alternative count. A combined source row can
be handled in at least three ways. It can remain combined, preserving the
source-observed table. It can be split for analytical purposes, producing a
derived view. Or it can be mistakenly treated as evidence that the source table
itself had separate f and v rows. The first two uses are legitimate under the
audit. The third is the drift error.

The hinge is attractive because it regularizes the table. A 27 by 8 view can
look more complete, more computationally convenient, and more compatible with
modern expectations about consonant rows. That attractiveness is exactly why it
needs a label. If a derived representation is cleaner than the source
representation, downstream systems will tend to prefer it. Without a layer
label, preference can become evidence. A teaching tool may present only the
cleaner view. A dataset may store only the derived rows. A later paper may cite
the derived total as though it were the original count.

The audit does not reject the f/v split. Rejecting every derived representation
would make useful analysis impossible. Instead, the audit asks the split to
carry its own conditions. The transformation is allowed when the combined row is
visible or recoverable, when the split is named, when the resulting count is
described as derived, and when no standards or source claim is inferred from the
derived count. A software system could satisfy this condition by storing a
``source row'' field, a ``derived row'' field, and a ``transformation'' field.
A manuscript can satisfy it by stating that 27/216 is an analytical expansion
from the combined f/v row.

This case study also shows why the audit belongs in digital humanities rather
than only in script encoding. The problem appears before formal encoding work
begins. It appears when a scholar writes a table caption, when a database schema
chooses a primary key, when an interface chooses a default view, and when a
paper compresses a layered evidence situation into one sentence. The f/v hinge
therefore functions as a small but concrete test of whether a representation
preserves provenance.

\section{Discussion}
The count-layer audit reframes the Nwagu Aneke count question. The issue is not
which number should win. The issue is which number answers which question. The
26 by 8 count answers the source-layer question under the present evidence
base. The 27/216 count answers a derived analytical question under the f/v
split. A future Unicode or standards count would answer a different question
and would require evidence that this paper does not claim to provide.

This reframing has practical consequences for scholarship. A source-critical
study should lead with the source-observed count. A standards-readiness study
should separate identity, repertoire, glyph evidence, names, usage, authority,
implementation, and proposal documentation. A computational study should state
whether it uses source or derived units. A cultural-heritage data model should
store the transformation rather than only the resulting integer. The same
source object can support all these studies only if its layers remain visible.

The audit also changes how negative results are treated. Blocking a 27/216
source claim is not a retreat from research. It is a positive result about the
limits of the current evidence. It prevents a downstream false premise. It also
makes future revision easier. If a later source review shows that separate f
and v rows are source-observed, the audit can be updated by changing the source
basis and revision history. If later review confirms that the split is derived,
the existing distinction becomes stronger. In both cases, the layer model gives
future evidence a place to land.

The broader implication is that cultural-heritage computing needs epistemic
typing. Metadata schemas and annotation models can store many things, but the
project must still decide whether an item is a source observation, editorial
normalization, analytical model, hypothesis, or blocked claim. When that
decision is left implicit, polished outputs can hide weak evidence. The risk is
heightened in AI-assisted writing and software generation, where systems can
produce coherent prose, tables, and interfaces from partially specified
materials. A clean output is not the same as a clean evidence chain.

\section{Implications for Standards and Encoding Discussion}
The audit should not be mistaken for a Unicode proposal. It is a pre-proposal
discipline for keeping the evidence legible. Unicode work requires character
identity, representative glyphs, names, properties, usage examples, stability,
implementation considerations, fonts, licensing, and proposal documentation
\cite{unicode17CoreSpec,unicodeProposals,unicodeRoadmaps,unicodePipeline2026,seiProposalTips}.
The present count-layer result supplies only one part of that future work: a
way to prevent a source table and a derived analytical view from being
confused.

That pre-proposal role is still useful. Without a count-layer audit, a future
standards discussion could begin from the wrong premise. It could mistake 216
derived cells for 216 observed source units. It could use a derived f/v split as
evidence for character separation before character-glyph analysis has been
completed. It could cite a source table while actually relying on an editorial
normalization. The audit prevents those moves by forcing every count to carry
its source basis and transformation rule.

The audit also helps identify the next standards-facing evidence. A future
readiness study would need source transcription review, public examples whose
rights are clear, representative glyph discussion, character naming principles,
directionality confirmation, punctuation and number evidence, usage examples,
community authority review, and implementation artifacts. Some of those items
are technical. Others are philological, legal, or social. Treating them as one
generic readiness problem would hide the work that remains.

\section{Implications for Digital Editions and Data Models}
A digital edition can implement the count-layer result directly. It can store
source rows and derived rows separately. It can show the 26 by 8 table as a
source view and the 27/216 f/v split as an analytical view. It can expose the
transformation rule and preserve links from derived cells back to source rows.
It can keep rights-sensitive source images or examples outside public release
until permissions and authority decisions are complete.

The same pattern applies to TEI, IIIF, Web Annotation, PROV-O, RO-Crate,
DataCite, and CIDOC CRM. TEI can distinguish between a source glyph occurrence,
an editorial label, and a derived representation. IIIF can identify source
regions when source images are available for the relevant use. Web Annotation
can attach a derived-cell assertion to a source target. PROV-O can record the
activity that split f/v. RO-Crate can package the audit with its source and
derived objects. DataCite can describe a citable release if one is approved.
CIDOC CRM can represent the cultural object, interpretation event, and actors.
The count-layer audit supplies the local epistemic decision these standards
would carry.

The practical design rule is simple: no interface should make the derived view
easier to see than the source relationship that produced it. If the 27/216 view
is displayed, the combined source row and the split operation should remain
recoverable. If a dataset exports derived rows, it should export the
transformation metadata. If a figure shows a normalized grid, its caption should
state that it is normalized. These design choices prevent visual regularity
from becoming false evidence.

\section{Implications for Computational Experiments}
Computational experiments often need clean units. A tokenizer baseline, search
index, OCR label set, or pedagogical interface may prefer the derived f/v view
because it gives a more regular structure. The audit does not forbid such use.
It requires the experiment to report what it used. A model trained on the source
view and a model trained on the derived view are not using the same evidence
layer. Their results should not be compared without naming that difference.

The count-layer distinction also affects evaluation. If a system emits 216 as
the source inventory, it has made a layer-promotion error. If it emits 216 as a
derived f/v analytical view and preserves the 26 by 8 source relationship, it
has preserved the evidence layer. A useful benchmark could therefore measure
promotion-error precision and recall, not just fluency or formatting. The
negative examples would include sentences and tables that look plausible while
promoting the derived count.

This point connects script documentation to broader concerns in language
technology. Unit choices shape downstream behavior. In subword tokenization,
the choice of segmentation affects model performance and representation
\cite{sennrich2016bpe,kudo2018sentencepiece,bostrom2020bytepair,rust2021tokenization}.
In citation-grounded generation, the presence of a citation does not guarantee
that the cited evidence supports the exact claim
\cite{wadden2020scifact,gao2023alce}. In an underdocumented script, the
presence of a clean count does not guarantee that the count belongs to the
source layer. The audit makes that distinction testable.

\section{Error Analysis}
The most likely error is not arithmetic. All relevant quantities are simple:
26 rows, 8 columns, 208 source-observed records, 27 derived rows, and 216 derived
cells. The difficult question is how those quantities travel through prose,
tables, filenames, interfaces, and later systems. Count-layer drift usually
appears when a correct number is moved into the wrong sentence.

The first error class is title and abstract compression. A careful methods
section may distinguish the source table from the derived f/v expansion, while
a title, abstract, or conclusion reports only the cleaner derived total. This is
dangerous because many readers encounter only the compressed surfaces of a
paper. A safe compressed statement says that the present source table has
208 records and that 216 is a derived f/v view. An unsafe compressed statement
says that the script has 216 signs. The latter may be shorter, but it changes
the evidence layer.

The second error class is table normalization. A normalized table can be easier
to read than the source table, especially when a derived f/v split makes the row
structure more regular. That visual improvement can mislead readers into
treating the normalized table as the observed one. The audit therefore requires
caption-level disclosure. A normalized table should state which source row was
split, what new rows were produced, and whether the source table remains
available for inspection. If the normalized view is the only view shown, the
paper should treat it as an analytical display rather than as a source record.

The third error class is data export loss. A spreadsheet or JSON export may
store derived rows without carrying the transformation metadata. Once the
transformation is lost, later tools have no way to know that 216 cells came from
a split. This error can persist even if the manuscript prose is careful. A data
model should therefore store at least four fields: source row identifier,
source-layer count, derived-row label, and transformation rule. If a data model
cannot store those fields, it should not be used as the authoritative carrier of
the result.

The fourth error class is citation displacement. A statement may cite Azuonye or
public script metadata while actually depending on a later analytical
transformation. The citation is real, but the claim is misplaced. This is the
same structural problem studied in claim-verification work: evidence can be
present while the exact claim outruns it \cite{wadden2020scifact,gao2023alce}.
For this paper, the remedy is to cite source identity sources for source
identity and transformation records for derived counts. A source citation alone
does not authorize the derived layer.

The fifth error class is standards leakage. Unicode and script-encoding
language is precise. Terms such as character, glyph, representative form,
property, repertoire, proposal, and implementation have consequences
\cite{unicode17CoreSpec,unicodeProposals}. A count-layer audit cannot turn a
source table into those objects. If standards language is used, it should be
used to identify missing work: character-glyph analysis, names, examples,
properties, rights, authority, and implementation evidence.

\section{Alternative Explanations}
One alternative explanation is that the difference between 208 and 216 is merely
a notation problem. On this view, a footnote could state that the numbers use
different conventions, and the rest of the paper could proceed with whichever
number is convenient. This paper rejects that approach because conventions
become claims when they enter titles, tables, datasets, and software. A
footnote is too weak if the exported artifact or default interface preserves
only the derived count. The layer must be represented structurally.

A second alternative explanation is that ordinary provenance metadata already
solves the problem. Provenance is necessary, but it is not sufficient. PROV-O
can record that a derived table was generated from a source table
\cite{w3cProvO}; RO-Crate can package both objects \cite{rocrate12}; TEI can
represent editorial choices \cite{teiP5Glyphs}. Those systems need the local
decision supplied by the audit: which count is observed, which count is
derived, and which public statements are permitted. Without that decision,
provenance machinery can faithfully preserve a mistaken claim.

A third alternative explanation is that the derived f/v split should simply be
discarded. That would be too conservative. Derived representations are often
valuable. They can support teaching, comparison, software tests, and
normalization. The problem is not derivation; the problem is unmarked
derivation. The audit preserves the useful derived view while preventing it
from becoming the source layer.

A fourth alternative explanation is that the evidence is too incomplete for any
paper. This objection is partly right and partly wrong. The evidence is too
incomplete for a repertoire proposal, public glyph corpus, or definitive
history of the script. It is not too incomplete for a source-critical method
paper about how counts should be typed. In fact, the incompleteness is why the
method is needed. A complete, rights-cleared, community-reviewed edition would
reduce the risk studied here; a partial evidence base makes the risk visible.

\section{Reviewer Objections}
A reviewer may object that the result is small. The response is that small
source-critical controls can have large downstream effects. If a derived count
is promoted early, later software, articles, and teaching materials may inherit
the mistake. Preventing one false foundation claim is a legitimate result when
the object is culturally significant and underdocumented.

A reviewer may object that the paper does not show enough primary source
material. That objection identifies a real limitation. The paper is not a
source-image edition. Its claim is about the relationship between a stated
source layer and a stated derived layer. If a future venue requires public
source images or glyph crops, rights review and source access must come first.
The present argument remains inspectable as a text-based layer audit, but it
does not substitute for image-based philology.

A reviewer may object that the f/v split requires linguistic or orthographic
analysis beyond the scope of a count audit. That is also fair. This paper does
not decide the full linguistic status of f and v in Nwagu Aneke. It decides the
status of the split as used in the current representation. Under the present
evidence base, the split is derived. A later linguistic or source review could
revise that status. The audit is designed so that such revision would change a
specific field rather than the whole paper silently.

A reviewer may object that digital-humanities standards already allow
uncertainty, so a new audit model is unnecessary. The response is that this
paper is not proposing a replacement standard. It is proposing a source-critical
decision that existing standards can carry. The standards answer how a decision
can be represented; the audit answers what the decision currently is.

\section{Future Evidence Work}
The next evidence task is source transcription review. A reviewer should inspect
the row structure, column structure, and combined f/v row against the source
dossier. The output should not be a polished paragraph. It should be a
disagreement matrix that records stable rows, uncertain rows, missing evidence,
and any row or column whose status changes under review.

The second task is rights and authority review. If future work needs to show
source images, glyph crops, or examples, it must establish what can be shown,
under what conditions, and with whose authority. A count-layer audit can guide
that work, but it cannot authorize public release. Rights and community
authority remain separate evidence categories.

The third task is structured implementation. The six-field audit model should be
encoded in a small machine-readable form: quantity, layer, source basis,
transformation rule, permitted wording, and revision condition. That form could
then be tested in TEI, JSON-LD, RO-Crate, or a relational schema. The test would
ask whether the model can export both source and derived views without losing
the transformation.

The fourth task is a promotion-error benchmark. Given a set of source-layer and
derived-layer statements, a system should identify whether a sentence preserves
or promotes the layer. The metric should report false negatives for dangerous
promotions and false positives for legitimate derived uses. Such a benchmark
would connect this source-critical paper to broader work on claim verification,
citation-grounded generation, and reliable computational representation.

\section{Limitations}
The audit is limited by the current evidence base. It accepts 26 by 8 as the
present source-observed layer for the purpose of this study, but future source
transcription review could revise rows, columns, readings, or the status of the
combined f/v row. The paper therefore does not claim a final critical edition
or a complete glyph corpus. It reports a layer decision under the evidence now
available.

The paper is also limited by its text-only scope. It does not reproduce source
images, glyph crops, manuscript examples, or a public dataset. Rights review
remains necessary before any public source-image or example release. Community
authority review remains necessary before stronger public representation claims
are made. These limitations are not peripheral. They define what kind of claim
the audit can responsibly support.

The six-field model is intentionally lightweight. It is sufficient for the
count problem studied here, but a full digital edition or standards dossier
would require more detailed metadata, source links, uncertainty measures,
rights statements, and authority records. A future version could formalize the
model in TEI, PROV-O, RO-Crate, or another structured format. The present paper
keeps the model small so that the source/derived distinction remains inspectable
without specialized software.

Finally, the audit does not decide whether the derived f/v split is
orthographically, pedagogically, or computationally superior. It only decides
how that split must be described if it is used. A future study could compare
teaching outcomes, search behavior, tokenizer coverage, or reader
interpretation under source and derived views. Such a study would need its own
data, methods, and ethics review.

\section{Reproducibility}
The reported result can be reproduced by checking five steps. First, identify
the current source-observed structure used by the study: 26 rows by
8 vowel/modifier columns, for 208 records. Second, identify the derived
operation: splitting a combined f/v row into separate analytical rows, yielding
27 rows and 216 cells. Third, inspect the manuscript for any sentence, caption,
or table that reports 27/216 as source-observed. Fourth, check that standards,
encoding, glyph, public-data, and authority claims are not inferred from either
count. Fifth, record the revision condition: if future source transcription
review changes the source table or the f/v interpretation, the count-layer
ledger and manuscript must change together.

A second reviewer can reproduce the standards-facing conclusion by comparing
the audit with Unicode and Script Encoding Working Group requirements
\cite{unicode17CoreSpec,unicodeProposals,unicodeRoadmaps,unicodePipeline2026,seiProposalTips}.
The count-layer audit does not supply representative glyphs, character names,
usage examples, implementation evidence, or a proposal dossier. It supplies a
source/derived boundary that such future work must preserve.

A third reviewer can reproduce the digital-humanities conclusion by asking how
the result would be represented in TEI, IIIF, Web Annotation, PROV-O, RO-Crate,
DataCite, or CIDOC CRM. In each case, the answer should preserve the
transformation from source layer to derived view rather than storing only a
single preferred count. If a representation cannot preserve that distinction,
it is not an adequate carrier for the result.

\section{Conclusion}
The Nwagu Aneke count problem is not resolved by choosing one attractive number.
It is resolved by assigning each number to the evidence layer it can support.
Under the current evidence base, the source-observed table is 26 by 8, yielding
208 records. The 27/216 representation is a derived f/v expansion. That
distinction lets digital-humanities, standards-readiness, and computational
work proceed without turning a useful analytical view into a false source
claim. The next scholarly work is source transcription review, rights review,
authority review, and more detailed representation of the ledger in a format
that preserves the same layer distinction.

\bibliographystyle{ACM-Reference-Format}
\bibliography{../references}
\end{document}
